\section{Frame and Model-Theoretic Semantics}
\label{sec:frame-model-semantics}
\slidesrc{Other semantics: 1--44; Semantics (long version): 7--10}

Distributional semantics (Section~\ref{sec:distributional-semantics}) tells us that
\emph{sold} is similar to \emph{bought}, but not \emph{who} sold \emph{what} to
\emph{whom}, nor whether a sentence is \emph{true}. This section presents two other views
of meaning, following slides 36--56 of Percy Liang's tutorial \emph{Natural Language
Understanding: Foundations and State-of-the-Art} (Simons Institute, 2017,
\url{https://simons.berkeley.edu/talks/percy-liang-01-27-2017-1}). The tutorial covers,
in order, properties of language, distributional semantics, frame semantics,
model-theoretic semantics, interactive learning and reflections; we use the parts on
frame and model-theoretic semantics.

This section covers:
\begin{enumerate}[nosep]
  \item \textbf{frame semantics} and semantic role labelling
        (Section~\ref{sec:frame-semantics});
  \item a comparison of the three types of semantics on one sentence
        (Section~\ref{sec:types-of-semantics});
  \item \textbf{model-theoretic semantics}: executable semantic parsing and
        compositionality (Section~\ref{sec:semantic-parsing}).
\end{enumerate}

\subsection{Frame Semantics}
\label{sec:frame-semantics}
\slidesrc{Other semantics: 2--18; Semantics (long version): 10}

\subsubsection{Word meaning revisited}
\slidesrc{Other semantics: 2--6}

What is the meaning of \emph{sold}?
\begin{itemize}
  \item \textbf{Distributional semantics}: all the contexts in which \emph{sold} occurs
        (\emph{\ldots was sold by\ldots}, \emph{\ldots sold me that piece of\ldots}). It can
        find similar words and contexts and generalise (e.g.\ with dimensionality
        reduction), but the representation is \textbf{monolithic}: a word vector has no
        internal structure, nothing in it says ``there is a seller and a buyer''.
  \item \textbf{Frame semantics}: the meaning is given by a \textbf{frame}, a
        stereotypical situation with \textbf{roles} (slots) to be filled.
\end{itemize}

\begin{nbdef}{Frame}
A \textbf{frame} is a structured representation of a stereotypical situation, made of
named \textbf{roles} (frame elements). A word \emph{evokes} the frame, and the sentence
fills its roles.
\end{nbdef}

\begin{nbexample}{a commercial transaction}
\emph{Cynthia sold the bike for \$200.}
\begin{center}
\begin{tabular}{@{}l l@{}}
\multicolumn{2}{@{}l}{\textbf{Commercial transaction}}\\
\midrule
\textsc{seller} & Cynthia \\
\textsc{buyer}  & ? \\
\textsc{goods}  & the bike \\
\textsc{price}  & \$200 \\
\end{tabular}
\end{center}
The \textsc{buyer} role stays empty: the sentence does not say who bought the bike, but
the frame tells us that \emph{someone} did.
\end{nbexample}

\subsubsection{Two properties of frames}
\slidesrc{Other semantics: 7--11}

Following Fillmore (1977) and Langacker (1987):

\begin{itemize}
  \item \textbf{Prototypical}: a frame does not need to handle all the cases. The frame
        of \emph{widow} is ``a woman marries one man, and the man dies''. What if a woman
        had three husbands, two of whom died? The frame does not say, and it does not
        need to: it describes the typical situation.
  \item \textbf{Profiling}: words evoking the same frame can \emph{highlight} different
        aspects of it.
        \begin{itemize}[nosep]
          \item \emph{sell} is seller-centric, \emph{buy} is buyer-centric:
                \emph{Cynthia sold the bike (to Bob)} vs.\ \emph{Bob bought the bike
                (from Cynthia)};
          \item \emph{rob} highlights the person, \emph{steal} the goods:
                \emph{Cynthia robbed Bob (of the bike)} vs.\ \emph{Cynthia stole the bike
                (from Bob)}.
        \end{itemize}
        The profiled role becomes the subject or object; the others are optional (in
        brackets).
\end{itemize}

\subsubsection{From syntax to semantics}
\slidesrc{Other semantics: 12--15}

The same situation, with \textsc{seller} Cynthia, \textsc{buyer} Bob, \textsc{goods} the
bike, \textsc{price} \$200, can be expressed by many \textbf{syntactic alternations},
with different verbs and different arguments in different positions:
\begin{itemize}[nosep]
  \item Cynthia sold the bike to Bob for \$200.
  \item The bike sold for \$200.
  \item Bob bought the bike from Cynthia.
  \item The bike was bought by Bob.
  \item The bike was bought for \$200.
  \item The bike was bought for \$200 by Bob.
\end{itemize}
The subject is the seller in the first sentence, the goods in the second, the buyer in
the third. \textbf{Goal}: map \emph{syntactic positions} to \emph{semantic roles}.

\subsubsection{Semantic role labelling}
\slidesrc{Other semantics: 16--18}

\begin{nbdef}{Semantic role labelling (SRL)}
Given a sentence, label each relevant phrase with its semantic role in the frame evoked
by the predicate:
\begin{center}
\begin{tabular}{@{}l ccccc@{}}
\textbf{Input:}  & Cynthia & sold & the bike & to Bob & for \$200 \\
\textbf{Output:} & \textsc{seller} & \textsc{predicate} & \textsc{goods} & \textsc{buyer} & \textsc{price}
\end{tabular}
\end{center}
Subtasks:
\begin{enumerate}[nosep]
  \item \textbf{frame identification}: find the \textsc{predicate} and the frame it
        evokes;
  \item \textbf{argument identification}: find the phrases filling \textsc{seller},
        \textsc{goods}, etc.
\end{enumerate}
\end{nbdef}

\begin{nbintuition}[SRL needs syntax]
\emph{Bob} is the \textsc{buyer} in ``Cynthia sold the bike \emph{to Bob}'' and in
``\emph{Bob} bought the bike'', but in different syntactic positions. SRL systems
therefore use the output of POS tagging and parsing
(Section~\ref{sec:sentence-level-syntax}): syntax is the input of semantics.
\end{nbintuition}

\textbf{Tools}: \textbf{FrameNet} (\url{https://framenet.icsi.berkeley.edu/}) is a
lexical database of English frames, with their roles and annotated example sentences;
see for example the frame \emph{Commerce\_buy}.

\subsection{Three Types of Semantics}
\label{sec:types-of-semantics}
\slidesrc{Other semantics: 19--22}

One sentence, three views:
\begin{center}
\emph{\textbf{Every} non-blue block is next to \textbf{some} blue block.}
\end{center}

\begin{itemize}
  \item \textbf{Distributional semantics}: \emph{block} is like \emph{brick},
        \emph{some} is like \emph{every}. True as far as usage goes (they appear in
        similar contexts), and fatal for meaning: swapping \emph{every} and \emph{some}
        changes what the sentence says completely.
  \item \textbf{Frame semantics}: \emph{is next to} has two arguments, \emph{block} and
        \emph{block}. It identifies the structure, but not the truth.
  \item \textbf{Model-theoretic semantics}: tells the difference between these two
        situations (``models''):
\end{itemize}

\begin{center}
\begin{tikzpicture}[blk/.style={draw, minimum size=0.7cm, font=\small\bfseries, text=white}]
  \node[blk, fill=nbred]   at (0,0)   {1};
  \node[blk, fill=nbblue]  at (0.75,0) {2};
  \node[blk, fill=nbgreen] at (2.0,0) {3};
  \node[blk, fill=nbblue]  at (2.75,0) {4};
  \node[font=\small] at (1.4,-0.8) {world A: the sentence is \textbf{true}};
  \begin{scope}[xshift=6cm]
  \node[blk, fill=nbred]   at (0,0)   {1};
  \node[blk, fill=nbred]   at (0.75,0) {2};
  \node[blk, fill=nbgreen] at (2.0,0) {3};
  \node[blk, fill=nbblue]  at (2.75,0) {4};
  \node[font=\small] at (1.4,-0.8) {world B: the sentence is \textbf{false}};
  \end{scope}
\end{tikzpicture}
\captionof{figure}{Two worlds; blocks 1--2 and 3--4 are next to each other.}
\end{center}

\begin{nbexample}{evaluating the sentence in each world}
In logic the sentence is
\[
  \forall x\, \bigl(\neg\mathit{blue}(x) \rightarrow \exists y\, (\mathit{blue}(y) \wedge
  \mathit{nextTo}(x, y))\bigr).
\]
The non-blue blocks are 1 and 3 in world A, 1, 2 and 3 in world B.
\begin{itemize}[nosep]
  \item World A: block 1 is next to 2 (blue), block 3 is next to 4 (blue): every
        non-blue block has a blue neighbour. \textbf{True}.
  \item World B: block 1 is next to 2, which is red, and to nothing else. \textbf{False}
        (one counterexample is enough for $\forall$).
\end{itemize}
\end{nbexample}

\begin{nbdef}{Model-theoretic semantics}
The meaning of a sentence is given by its \textbf{truth conditions}: the set of
situations (\emph{models} of the world) in which it is true. Understanding a sentence
means being able to decide, given a model, whether it is true; more generally, to map
the sentence to a formal representation (a logical form, a query, a program) that can be
\emph{evaluated} against the world.
\end{nbdef}

\begin{center}
\small
\begin{tabularx}{\linewidth}{@{}l >{\raggedright\arraybackslash}X >{\raggedright\arraybackslash}X@{}}
\toprule
\textbf{Type} & \textbf{Represents a meaning as\ldots} & \textbf{Good at / misses} \\
\midrule
Distributional & a vector of contexts (Section~\ref{sec:distributional-semantics}) &
  similarity, learnt automatically / no structure, no truth \\
Frame & a frame with roles (Section~\ref{sec:frame-semantics}) &
  who did what to whom / no logic (\emph{every}, \emph{not}) \\
Model-theoretic & truth conditions, a logical form or program &
  precise reasoning, answers / hard to obtain from free text \\
\bottomrule
\end{tabularx}
\captionof{table}{The three types of semantics compared.}
\end{center}

\subsection{Model-Theoretic Semantics in Practice: Semantic Parsing}
\label{sec:semantic-parsing}
\slidesrc{Other semantics: 23--44; Semantics (long version): 10}

\subsubsection{Executable semantic parsing}
\slidesrc{Other semantics: 23--34}

\begin{nbdef}{Executable semantic parsing}
\textbf{Semantic parsing} maps a natural-language sentence to a \textbf{program}
(logical form) in a formal language. The program is then \textbf{executed} in a
context (a database, a calendar, a robot's world) to produce a behaviour: an answer, an
action.
\[
  \text{[context]} \;+\; \text{[sentence]}
  \;\xrightarrow{\text{semantic parsing}}\; \text{[program]}
  \;\xrightarrow{\text{execute}}\; \text{[behaviour]}
\]
\end{nbdef}

\begin{nbexample}{a question to a database}
\emph{What is the largest city in Europe by population?} The program is built piece by
piece:
\begin{itemize}[nosep]
  \item $\textit{Cities}$: the set of all cities;
  \item $\textit{ContainedBy}(\textit{Europe})$: the set of things located in Europe;
  \item $\textit{Cities} \cap \textit{ContainedBy}(\textit{Europe})$: cities in Europe;
  \item the one with the largest population:
\end{itemize}
\[
  \argmax\bigl(\textit{Cities} \cap \textit{ContainedBy}(\textit{Europe}),\ \textit{Population}\bigr).
\]
Executed on a geographical database it returns \textbf{Istanbul}.
\end{nbexample}

\begin{nbexample}{a command to a calendar}
\emph{Remind me to buy milk after my last meeting on Monday.}
\[
  \textit{Add}\bigl(\textit{Buy}(\textit{Milk}),\ \argmax(\textit{Meetings} \cap
  \textit{HasDate}(\text{2016-07-18}),\ \textit{EndTime})\bigr)
\]
``My last meeting on Monday'' becomes the meeting of that date with the largest end
time; executing the program adds the reminder.
\end{nbexample}

Note the context: ``Monday'' became a precise date (2016-07-18) only because the system
knows today's date; ``my meetings'' are the ones in \emph{this} user's calendar.

\begin{nbwhy}
This is exactly what ChatBDI does (Section~\ref{sec:chatbdi}): it turns a sentence into a
term such as \texttt{tire\_pressure(front, \_)} that an agent can execute. It is also how
voice assistants and ``text-to-SQL'' systems work, and how LLMs call tools (they generate
a function call, which is then executed).
\end{nbwhy}

\subsubsection{A brief history of semantic parsing}
\slidesrc{Other semantics: 35}

\begin{center}
\small
\begin{tabularx}{\linewidth}{@{}>{\raggedright\arraybackslash}X >{\raggedright\arraybackslash}X@{}}
\toprule
GeoQuery & Zelle \& Mooney 1996 \\
Inductive logic programming & Tang \& Mooney 2001 \\
CCG & Zettlemoyer \& Collins 2005 \\
String kernels & Kate \& Mooney 2006 \\
Synchronous grammars; relaxed CCG & Wong \& Mooney 2007; Zettlemoyer \& Collins 2007 \\
Learning from the world & Clarke et al.\ 2010 \\
Higher-order unification; learning from answers & Kwiatkowski et al.\ 2011; Liang et al.\ 2011 \\
Language + vision & Matuszek et al.\ 2012 \\
Large-scale KBs; regular expressions; instruction following & 2013 (Berant et al., Kwiatkowski et al.; Kushman et al.; Artzi \& Zettlemoyer) \\
Reduction to paraphrasing & Berant \& Liang 2014 \\
Compositionality on tables; dataset from logical forms & Pasupat \& Liang 2015; Wang et al.\ 2015 \\
\bottomrule
\end{tabularx}
\end{center}

The trend: from small hand-annotated datasets (GeoQuery: questions about US geography,
each paired with its logical form) to learning from weaker supervision (only the answers,
or the effect in the world) and to large knowledge bases.

\subsubsection{Compositional semantics}
\slidesrc{Other semantics: 36--41}

How do we get from the sentence to the program? The principle of
\textbf{compositionality}, formalised by \textbf{Richard Montague}: the meaning of a
phrase is a function of the meanings of its parts and of the way they are combined
syntactically.

\begin{nbexample}{cities in Europe}
\begin{center}
\begin{tikzpicture}[level distance=1.1cm,
    level 1/.style={sibling distance=4.4cm},
    level 2/.style={sibling distance=2.6cm},
    every node/.style={font=\small, align=center}]
  \node {$\textit{Cities} \cap \textit{ContainedBy}(\textit{Europe})$}
    child {node {$\textit{Cities}$\\[-2pt]{\color{nbgray}\emph{cities}}}}
    child {node {$\textit{ContainedBy}(\textit{Europe})$}
      child {node {$\textit{ContainedBy}$\\[-2pt]{\color{nbgray}\emph{in}}}}
      child {node {$\textit{Europe}$\\[-2pt]{\color{nbgray}\emph{Europe}}}}};
\end{tikzpicture}
\captionof{figure}{Bottom-up composition: each word gets a meaning from the lexicon
(\emph{cities} $\mapsto$ \textit{Cities}, \emph{in} $\mapsto$ \textit{ContainedBy},
\emph{Europe} $\mapsto$ \textit{Europe}), and each node of the tree combines the meanings
of its children (function application, then intersection).}
\end{center}
\end{nbexample}

The tree is the parse tree of Section~\ref{sec:grammars-parsing}, decorated with
meanings: once again, syntax is the scaffolding of semantics.

\subsubsection{Language variation}
\slidesrc{Other semantics: 42--44}

The same meaning,
\[ \textit{Cities} \cap \textit{ContainedBy}(\textit{Europe}), \]
can be expressed in many ways:
\begin{itemize}[nosep]
  \item cities in Europe
  \item European cities
  \item cities that are in Europe
  \item cities located in Europe
  \item cities on the European continent
\end{itemize}
Can we write composition rules for all of these by hand? The slides answer:
\textbf{no!} \emph{European} is a single word that must produce
$\textit{ContainedBy}(\textit{Europe})$; \emph{located in} and \emph{on the \ldots\
continent} are other ways of saying \emph{in}; and natural language has endlessly many
such variants. This is why modern semantic parsers \emph{learn} the mapping from data
(and today LLMs do it from examples in the prompt), keeping the logical form as the
target.

\textbf{Tools}: for model-theoretic semantics the slides report ``\ldots not found\ldots'':
there is no standard off-the-shelf tool comparable to FrameNet.

\begin{nbsummary}[frame and model-theoretic semantics]
\begin{itemize}[nosep]
  \item Distributional vectors are monolithic. \textbf{Frame semantics}: a word evokes a
        frame (stereotypical situation) with roles (\textsc{seller}, \textsc{buyer},
        \textsc{goods}, \textsc{price}). Frames are \textbf{prototypical} (\emph{widow})
        and words \textbf{profile} roles (\emph{sell}/\emph{buy}, \emph{rob}/\emph{steal}).
  \item Many syntactic alternations, one frame: \textbf{semantic role labelling} maps
        syntactic positions to roles (frame identification + argument identification).
        Tool: FrameNet.
  \item ``Every non-blue block is next to some blue block'': distributional (block
        $\approx$ brick), frame (\emph{next to}(block, block)), model-theoretic (true or
        false in a given world).
  \item \textbf{Model-theoretic semantics}: meaning = truth conditions. Executable
        semantic parsing: sentence $\to$ program $\to$ execution (database, calendar).
  \item \textbf{Compositionality} (Montague): meaning of the whole from the meanings of
        the parts along the parse tree. Language variation makes hand-written rules
        impossible: the mapping is learnt.
\end{itemize}
\end{nbsummary}

